% =============================================================
% SECCIÓN: Idiomas y Habilidades Destacadas
% Archivo: sections/otros/5_idiomas_habilidades.tex
% Descripción: Niveles de idioma y competencias blandas relevantes
%              según el perfil de postulación.
% Escala de skills: 1 (básico) → 5 (nativo/experto)
% Actualizado: 2026
% =============================================================

\twocolumnsection
{\sectionTitle{Idiomas}{\faLanguage}
	% ----------------------------------------------------------
	% COLUMNA IZQUIERDA: Idiomas
	% ----------------------------------------------------------
	\begin{skills}
		% Quechua — lengua materna de la región
		\skill{Quechua}{5}
		% Español — nativo
		\skill{Español}{5}
		% Inglés — nivel básico-intermedio (ICPNA hasta B10, 2019-2022)
		\skill{Inglés}{3}
	\end{skills}}
{\sectionTitle{Habilidades Destacadas}{\faPlus}
	% ----------------------------------------------------------
	% COLUMNA DERECHA: Habilidades Destacadas
	% Adapta la lista según el perfil de postulación.
	% ----------------------------------------------------------
	\vspace{1em}
	\begin{itemize}
		% Habilidades transversales (aplican a todos los perfiles)
		\item Capacidad analítica
		\item Iniciativa en proyectos de datos
		\item Trabajo en equipo
		\item Adaptación rápida a nuevas herramientas
		\item Organización al detalle en bases de datos y reportes
	\end{itemize}
}



		%----------------------------------------------------------
		% Habilidades específicas para DOCENCIA (descomentar si aplica)
		%----------------------------------------------------------
		%\item Planificación y ejecución de sesiones de aprendizaje
		%\item Manejo de plataformas virtuales (Moodle, Google Classroom)
		%\item Uso de Inteligencia Artificial como herramienta pedagógica
		
		%----------------------------------------------------------
		% Habilidades específicas para DATA SCIENCE (descomentar si aplica)
		%----------------------------------------------------------
		%\item Programación en R y Python para análisis de datos
		%\item Visualización de datos con ggplot2 y herramientas afines
		%\item Procesamiento de microdatos (ENAHO, INEI, BCRP)